% ============================================================
%  Chapter 2: Literature Review
%  Henry Kabuye PhD Thesis — Northumbria University
%  AI-Based Intelligent Multi-Stage and User-Centric
%  Ransomware Attack Detection
% ============================================================

\chapter{Literature Review}
\label{chap:litreview}

\lettrine{R}{ansomware} has emerged as one of the most operationally disruptive and economically
damaging categories of malicious software confronting organisations, governments, and critical
infrastructure operators worldwide.
Unlike earlier forms of financial malware that sought to exfiltrate credentials or conduct
fraudulent transactions silently, ransomware explicitly attacks data availability: it encrypts
or exfiltrates victim files and demands a ransom payment in exchange for the decryption key or
the promise of non-publication \cite{beaman2021ransomware, greenberg2018untold}.
The detection challenge it poses is correspondingly severe: the encryption payload must be
identified and contained within a narrow window before irrecoverable damage is done, yet the
early-stage behavioural and structural signals that precede encryption overlap substantially
with those of legitimate software, making high-precision early detection technically demanding.

This chapter surveys the body of knowledge on which the empirical investigations in this
thesis are built.
Section~\ref{sec:lit_threat} characterises the ransomware threat landscape — its taxonomy,
attack lifecycle, and documented impact.
Section~\ref{sec:lit_ml} reviews machine learning approaches to malware and ransomware
detection from classical methods through to contemporary deep learning architectures.
Section~\ref{sec:lit_uq} covers the theory and practice of uncertainty quantification in
deep neural networks, with particular attention to Monte Carlo Dropout and probability
calibration.
Section~\ref{sec:lit_xai} addresses explainable artificial intelligence for cybersecurity,
covering feature attribution, counterfactual methods, and their application to malware
analysis.
Section~\ref{sec:lit_mas} reviews multi-agent architectures and modality fusion strategies
relevant to multi-signal detection.
Section~\ref{sec:lit_kbs} surveys neuro-symbolic AI and the operationalisation of
threat-intelligence knowledge as a detection signal.
Section~\ref{sec:lit_cd} reviews concept drift and continual learning in the context of
evolving malware populations.
Section~\ref{sec:lit_gaps} synthesises the identified gaps and maps them to the three
empirical contributions of this thesis.

% ============================================================
\section{The Ransomware Threat Landscape}
\label{sec:lit_threat}
% ============================================================

% ------------------------------------------------------------
\subsection{Taxonomy and Attack Lifecycle}
\label{sec:lit_taxonomy}
% ------------------------------------------------------------

The term \emph{ransomware} encompasses a heterogeneous family of malicious programs united by
the common goal of extorting payment from victims in exchange for restoring access to their
own data or systems \cite{razaulla2023age}.
Early academic taxonomies distinguished two primary categories: \emph{locker ransomware}, which
denies access to the system interface or desktop without modifying file contents, and
\emph{crypto-ransomware}, which selectively or exhaustively encrypts user files using
asymmetric or hybrid cryptographic schemes \cite{moussaileb2021survey, sikorski2012practical}.
Contemporary ransomware campaigns have largely converged on the crypto variant, augmented by
\emph{double-extortion} tactics in which data is exfiltrated prior to encryption so that
public release of sensitive materials constitutes a second leverage mechanism independent of
the decryption key \cite{beaman2021ransomware, meland2020ransomware}.

The attack lifecycle of a ransomware campaign maps closely to the MITRE ATT\&CK Enterprise
framework \cite{razaulla2023age}.
An attack typically begins with initial access via phishing, remote desktop protocol
exploitation, or compromised credentials, followed by execution of a dropper or loader
component.
The ransomware then performs reconnaissance — enumerating network shares, user directories,
and backup locations — before moving laterally to maximise the encrypted surface area.
Privilege escalation and defence evasion tactics are employed to disable endpoint protection
and shadow copy services.
Finally, the encryption payload is deployed and a ransom note delivered.
The window between initial execution and the completion of encryption can be as short as
minutes for opportunistic campaigns, making real-time or near-real-time detection
essential \cite{SgandurraMML16, hirano2022ransap}.

A number of taxonomies have been proposed specifically to characterise the technical
diversity of ransomware families.
Sgandurra et al.\ \cite{SgandurraMML16} distinguish families by the cryptographic scheme
employed (symmetric, asymmetric, or hybrid), the file targeting strategy (by extension, by
directory, or exhaustive), and the command-and-control communication protocol.
More recent surveys \cite{yumlembam2024comprehensive, Ozkan-Okay2024Survey} extend these
taxonomies to cover ransomware-as-a-service (RaaS) distribution models, which have
substantially lowered the technical barrier to launching campaigns and accelerated the
proliferation of novel variants.
Polymorphic and metamorphic ransomware families that mutate their code at each deployment to
evade signature-based detection represent a particularly acute challenge for static analysis
systems \cite{Urooj2024polymorphism}.

% ------------------------------------------------------------
\subsection{Economic and Operational Impact}
\label{sec:lit_impact}
% ------------------------------------------------------------

The economic cost of ransomware is substantial and escalating.
Europol's Internet Organised Crime Threat Assessment documents ransomware as the dominant
cyber-extortion threat facing European organisations, with critical infrastructure sectors
including healthcare, energy, and transportation identified as priority targets
\cite{europol2020iocta}.
The 2017 WannaCry campaign — estimated to have infected over 300,000 systems across 150
countries in a single day — caused disruption to the United Kingdom's National Health Service
that resulted in the cancellation of nearly 20,000 outpatient appointments and cost an
estimated \pounds 92 million in remediation and lost output \cite{greenberg2018untold}.
The 2021 Colonial Pipeline attack demonstrated that ransomware could directly impact physical
infrastructure, forcing a major fuel supply pipeline operator to halt operations for nearly a
week and triggering regional fuel shortages across the eastern United States \cite{cdk_global_attack}.

Beyond direct financial losses, ransomware imposes substantial indirect costs through
reputational damage, regulatory fines for inadequate data protection, increased insurance
premiums, and the allocation of significant analyst time to incident response and forensic
investigation \cite{beaman2021ransomware, europol2020iocta}.
The operational asymmetry is stark: a successful ransomware campaign may be executed with
modest technical resources by a RaaS affiliate, yet its remediation may require weeks of
specialist effort and expenditure far exceeding the initial ransom demand.
This asymmetry directly motivates the development of automated, high-throughput detection
frameworks capable of operating with minimal analyst involvement under normal conditions
while directing human attention precisely toward the cases in which it is most needed.

% ============================================================
\section{Machine Learning for Malware and Ransomware Detection}
\label{sec:lit_ml}
% ============================================================

% ------------------------------------------------------------
\subsection{Classical Machine Learning Approaches}
\label{sec:lit_classical}
% ------------------------------------------------------------

The application of machine learning to malware detection predates the deep learning era by
over a decade.
Early systems employed classical discriminative models — decision trees, na\"{i}ve Bayes
classifiers, support vector machines, and random forests — trained on handcrafted feature
vectors derived from static program analysis \cite{sarker2020cybersecurity}.
Static features commonly used in this period included import address table n-gram profiles,
section entropy values, printable string statistics, and PE header field values; dynamic
features derived from sandbox execution logs included system call sequences, registry
modification patterns, and network traffic summaries \cite{sikorski2012practical}.

Random forests proved particularly effective for static malware classification due to their
robustness to high-dimensional sparse feature spaces and their natural support for feature
importance estimation \cite{fernando2020study}.
Gradient-boosted ensembles such as XGBoost \cite{chen2016xgboost} extended this performance
further on structured tabular data by capturing non-linear interactions between features
through sequential tree residual fitting.
Studies specifically targeting ransomware classification demonstrated that dynamic behavioural
features — in particular, file system access patterns and entropy measurements on written
data — are strongly discriminative of ransomware activity when evaluated against legitimate
software \cite{SgandurraMML16, krishna2020comparative}.

However, classical machine learning approaches exhibit two well-documented limitations in the
ransomware context.
First, they require extensive feature engineering expertise to achieve competitive performance:
effective feature selection depends on domain knowledge about which program behaviours are
indicative, and this knowledge must be continuously updated as adversaries adapt their
tactics \cite{razaulla2023age, gibert2020rise}.
Second, classical models are particularly susceptible to concept drift: as ransomware families
evolve and novel variants are deployed, the feature distribution of the malicious class shifts
away from the distribution seen during training, and model accuracy degrades
measurably \cite{Ceschin2022ConceptDrift, Fernando2024FeSAD}.
These limitations motivate the transition to representation-learning approaches that can
acquire task-relevant features automatically and be updated as the threat landscape evolves.

% ------------------------------------------------------------
\subsection{Deep Learning Approaches}
\label{sec:lit_deep}
% ------------------------------------------------------------

Deep neural networks offer two principal advantages over classical models for malware
detection.
First, they can learn hierarchical feature representations directly from raw or minimally
preprocessed inputs, reducing dependence on expert feature engineering \cite{gibert2020rise}.
Second, architectures such as convolutional neural networks (CNNs) and recurrent networks
can exploit spatial and temporal structure in the input data that is ignored by conventional
flat feature vectors.

CNNs have been applied to malware classification by treating the byte stream of an executable
file as a one-dimensional or two-dimensional signal and learning local pattern detectors via
convolutional filters \cite{kolosnjaji2016deep}.
This approach achieves state-of-the-art accuracy on family classification benchmarks without
any hand-specified features, and has been extended to ransomware-specific tasks with strong
results \cite{zahoora2022Deep}.
Recurrent architectures, and particularly Long Short-Term Memory (LSTM) networks, have been
applied to system call sequence data to exploit the temporal ordering of API calls during
execution — capturing the characteristic enumeration-then-encryption sequence that precedes
ransomware payload delivery \cite{hirano2022ransap, HiranoK22}.

More recently, attention mechanisms and transformer-based architectures \cite{vaswani2017attention}
have been explored for security log analysis and network intrusion detection, leveraging
self-attention to model long-range dependencies in event sequences that may span hours of
activity before an attack payload executes \cite{tian2025llmsurvey, ZhangLZYHY25}.
LLM-derived embeddings of log data have also been investigated as input representations for
downstream classifiers \cite{tian2025llmsurvey}, though these systems are predominantly
evaluated as standalone classifiers without uncertainty-aware fusion or structured triage
logic.

Despite strong discriminative performance, deep learning classifiers for ransomware detection
share a critical limitation: they produce overconfident, poorly calibrated posterior
probability estimates, particularly for inputs near or between the training class
distributions \cite{gawlikowski2023Uncertainity, guo2017calibration}.
A model that assigns 99\% confidence to a borderline prediction provides the same output
format as one making a prediction on a clear-cut case, giving the analyst no signal about
when autonomous decisions are reliable and when human review should be sought.
This limitation is the primary motivation for the uncertainty quantification methods reviewed
in Section~\ref{sec:lit_uq}.

% ------------------------------------------------------------
\subsection{Feature Representations for Ransomware}
\label{sec:lit_features}
% ------------------------------------------------------------

A consistent finding across both classical and deep learning ransomware detection studies is
that no single feature modality captures all discriminative signal, and that different
modalities are most informative at different stages of the attack lifecycle.

\emph{Static} or \emph{semantic} features, derived from analysis of the executable without
running it, provide rapid pre-execution screening but are susceptible to obfuscation: packing,
encryption, and code mutation can substantially alter the static feature profile of a
ransomware binary without changing its runtime behaviour \cite{Urooj2024polymorphism,
sikorski2012practical}.
Studies using PE-format features — section entropy, import table statistics, and PE header
fields — report high detection rates on known families but degraded performance on novel or
obfuscated samples \cite{GogineniDV22ForeseerRanSAP}.

\emph{Dynamic} or \emph{behavioural} features, derived from monitoring the execution of the
binary in a controlled environment, are considerably more robust to static obfuscation but
introduce a detection latency: the sample must be executed, at least partially, before
behavioural features can be collected \cite{hirano2022ransap}.
System API call traces, file access patterns, network connection sequences, and registry
modification logs all carry strong ransomware-indicative signal, with file I/O entropy
elevation being particularly reliable as an indicator of encryption activity
\cite{SgandurraMML16, hirano2022ransap}.

A substantial literature demonstrates that \emph{hybrid} approaches that combine static and
dynamic features outperform either modality
alone \cite{Ferdous2024Comp, yumlembam2024comprehensive}.
However, most hybrid approaches in the literature combine features from both modalities into
a single flat vector fed to a single unified model, losing the modality-specific
distributional characteristics and preventing the system from reasoning separately about
confidence in the static and dynamic signals.
The multi-agent architecture introduced in this thesis addresses this limitation by training
independent agents on each modality and fusing their outputs at the score level, preserving
modality-specific uncertainty estimates that are used to govern triage decisions.

% ============================================================
\section{Uncertainty Quantification in Deep Neural Networks}
\label{sec:lit_uq}
% ============================================================

% ------------------------------------------------------------
\subsection{Sources and Types of Uncertainty}
\label{sec:lit_uq_types}
% ------------------------------------------------------------

The uncertainty in the predictions of a deep neural network arises from two conceptually
distinct sources \cite{gawlikowski2023Uncertainity, malinin2018predictive}.
\emph{Aleatoric} uncertainty, also called data or irreducible uncertainty, reflects genuine
ambiguity in the input: when a sample lies near or on the decision boundary between classes,
the labelling itself is inherently uncertain and no model can resolve it without additional
information.
\emph{Epistemic} uncertainty, also called model or reducible uncertainty, reflects the
model's lack of knowledge: it arises when a sample falls outside or at the periphery of the
training distribution, so the model has insufficient evidence to form a reliable
estimate \cite{gawlikowski2023Uncertainity}.

In the ransomware detection context, both types of uncertainty carry operational significance.
Aleatoric uncertainty indicates borderline samples that may require additional dynamic
analysis or analyst review regardless of how much training data is available.
Epistemic uncertainty indicates novel or out-of-distribution samples — potentially
representing a new ransomware family or a benign program that has not previously been
seen — where the model's prediction should be treated with
scepticism \cite{Doan2024uncertainBay}.
A detection framework that can distinguish between these two types of uncertainty provides
more actionable intelligence than one that reports only a single scalar confidence score.

% ------------------------------------------------------------
\subsection{Monte Carlo Dropout}
\label{sec:lit_mcd}
% ------------------------------------------------------------

Gal and Ghahramani \cite{gal2016dropout} demonstrated that a neural network trained with
dropout regularisation can be interpreted as an approximate Bayesian neural network: if
dropout masks are sampled stochastically at inference time rather than fixed or averaged as
in standard practice, the resulting distribution over predictions approximates the posterior
predictive distribution of a Bayesian model over network weights.
This insight, known as Monte Carlo (MC) Dropout, enables epistemic uncertainty estimation
from any standard dropout-regularised network without modification to the training procedure
and with minimal computational overhead.

At inference time, $T$ forward passes are performed through the network with independently
sampled dropout masks, yielding $T$ class probability estimates $\hat{p}_1, \ldots, \hat{p}_T$.
The predictive mean $\bar{p} = \frac{1}{T}\sum_{t=1}^{T}\hat{p}_t$ serves as the point
prediction, while the predictive standard deviation
$\sigma = \sqrt{\frac{1}{T}\sum_{t=1}^{T}(\hat{p}_t - \bar{p})^2}$ quantifies the
epistemic uncertainty of that prediction.
High values of $\sigma$ on a given input indicate that the dropout ensemble disagrees
substantially about the correct prediction, which is characteristic of out-of-distribution
or borderline inputs \cite{gal2016dropout, gawlikowski2023Uncertainity}.

MC Dropout has been applied to a range of security tasks including anomaly detection,
network intrusion classification, and, more recently, malware
family identification \cite{Doan2024uncertainBay}.
The number of Monte Carlo samples $T$ governs the trade-off between uncertainty estimation
quality and inference latency: empirical studies suggest that $T = 20$--$50$ passes are
sufficient for stable uncertainty estimates in most classification settings, with diminishing
returns beyond this range \cite{gal2016dropout}.
The method is particularly well suited to the multi-agent framework developed in this thesis,
where uncertainty estimates produced by individual agents are used both for triage decisions
and for inter-agent communication through the Peer Consultation Protocol.

% ------------------------------------------------------------
\subsection{Probability Calibration}
\label{sec:lit_calibration}
% ------------------------------------------------------------

A well-calibrated classifier is one whose posterior probability estimates are consistent with
empirical frequencies: a batch of inputs to which the model assigns probability 0.8 should,
on average, contain approximately 80\% positives \cite{guo2017calibration}.
The majority of deep neural networks trained with maximum likelihood objectives are
poorly calibrated, exhibiting a systematic tendency toward overconfidence: the model assigns
posterior probabilities close to 0 or 1 far more frequently than the actual frequency of
correct predictions in those ranges would
justify \cite{guo2017calibration, malinin2018predictive}.
This overconfidence is exacerbated in the presence of class imbalance — a common condition
in ransomware datasets — because the model is incentivised to minimise the cross-entropy loss
on the dominant class at the expense of calibration on the minority class.

Guo et al.\ \cite{guo2017calibration} demonstrated that \emph{temperature scaling}, a
post-hoc calibration method that divides the model's logit output by a scalar temperature
parameter $T_{\mathrm{cal}}$ fitted on a held-out validation set, substantially reduces
calibration error across a range of deep learning architectures with negligible impact on
discriminative accuracy.
The method is particularly attractive for multi-agent systems in which the probability
outputs of independently trained agents must be combined, because it ensures that the fusion
layer receives probability estimates that are comparable in their interpretive meaning across
agents rather than reflecting idiosyncratic calibration errors specific to each agent's
training regime.

% ============================================================
\section{Explainable Artificial Intelligence for Cybersecurity}
\label{sec:lit_xai}
% ============================================================

% ------------------------------------------------------------
\subsection{Feature Attribution Methods}
\label{sec:lit_shap}
% ------------------------------------------------------------

The opacity of deep neural network predictions — the so-called ``black-box'' property — is a
significant practical barrier to the deployment of neural network-based malware detectors in
operational security environments.
Security analysts making containment or escalation decisions require not only a prediction
but an explanation of the evidence underlying it: which features drove the ransomware
classification, and how much would their values need to change before the classification
reversed \cite{lundberg2017shap, ribeiro2016lime}?
Explainable AI (XAI) methods address this requirement by associating a predicted label with
an interpretable description of its drivers.

SHAP (SHapley Additive exPlanations) \cite{lundberg2017shap} grounds attribution in
cooperative game theory: each feature's contribution to a prediction is defined as the
Shapley value — the weighted average marginal contribution of that feature across all
possible feature subsets.
This attribution is guaranteed to satisfy a set of desirable axiomatic properties (efficiency,
symmetry, dummy, and linearity) that make the resulting explanations consistent and
non-arbitrary.
TreeSHAP, DeepSHAP, and KernelSHAP variants provide efficient computation for tree ensembles,
deep networks, and model-agnostic settings respectively.
For ransomware detection, SHAP feature importance plots can identify which PE-header fields or
API call patterns are driving a detection, providing actionable intelligence for signature
refinement and analyst triage.

LIME (Local Interpretable Model-agnostic Explanations) \cite{ribeiro2016lime} takes a
complementary approach: for a specific input, LIME perturbs the input in the neighbourhood
of its feature values, queries the black-box model on the perturbed samples, and fits a
locally linear interpretable model to the resulting predictions.
The coefficients of this local linear model constitute the explanation for that specific
prediction.
LIME is more computationally expensive than SHAP for repeated predictions but offers the
advantage of being fully model-agnostic and producing explanations that are directly
readable as weighted feature lists.

Both SHAP and LIME explanations have been applied to malware and ransomware detection tasks
to validate model behaviour, identify spurious correlations in training data, and generate
feature-level intelligence for analyst use \cite{gulmez2024xran, Ispahany2024Survey}.

% ------------------------------------------------------------
\subsection{Counterfactual Explanations}
\label{sec:lit_cf}
% ------------------------------------------------------------

While feature attribution methods identify \emph{which} features drove a prediction, they do
not directly answer the analyst's operational question: \emph{what would need to be different
about this sample for it to be classified as benign?}
Counterfactual explanations address this question by finding the minimal perturbation of the
input feature vector that would reverse the model's prediction, subject to the constraint that
the perturbed input remains plausible given the known distribution of benign
samples \cite{Farzadpour2024Thresholding}.

Counterfactual explanations are particularly valuable in the ransomware context because they
translate model behaviour into analyst-actionable intelligence.
A counterfactual that identifies two specific PE-header flag values as the minimal change
needed to reclassify a suspicious executable as benign gives an analyst a targeted hypothesis
to investigate; a counterfactual that requires altering 40 features by large amounts suggests
that the sample is confidently ransomware and can be contained with high precision.
The distance from an input to its nearest counterfactual (typically measured in $L_2$ or
$L_0$ norm) also provides an interpretable confidence proxy: inputs with close counterfactuals
lie near the decision boundary and warrant greater analyst scrutiny.

% ------------------------------------------------------------
\subsection{Explainability in Malware Analysis}
\label{sec:lit_xai_malware}
% ------------------------------------------------------------

The application of XAI methods to malware analysis has grown substantially since 2020,
driven by regulatory expectations around AI transparency and the practical demand for
analyst-comprehensible outputs from automated security systems.
Recent surveys \cite{Ispahany2024Survey, yumlembam2024comprehensive} identify feature
attribution-based explanation as the most widely adopted approach in the malware detection
literature, followed by rule extraction, attention visualisation, and counterfactual generation.
XRan \cite{gulmez2024xran} specifically applies SHAP to explain ransomware classifications
derived from API call profiles, demonstrating that attribution values correlate with known
ransomware-indicative API families and can be used to cluster samples by behavioural similarity.

A consistent finding in this literature is that explanations produced by attribution methods
are most useful when they are grounded in the analyst's existing knowledge framework: SHAP
values over raw byte offsets are less actionable than SHAP values over named PE-format fields,
even if both represent the same underlying model behaviour \cite{lundberg2017shap, gulmez2024xran}.
This observation motivates the feature-level explanation design choices in Chapter~3 of this
thesis, where explanations are presented over named API call groups and PE-format categories
rather than raw feature indices.

% ============================================================
\section{Multi-Agent Systems and Modality Fusion}
\label{sec:lit_mas}
% ============================================================

% ------------------------------------------------------------
\subsection{Multi-Agent Architectures for Detection}
\label{sec:lit_mas_arch}
% ------------------------------------------------------------

Multi-agent systems distribute a complex detection or classification task across a collection
of autonomous agents, each responsible for a subset of the problem, with coordination
mechanisms determining how individual agent outputs are combined into a system-level
decision \cite{MultiAgentAdaptiveDetection2023}.
In the cybersecurity domain, multi-agent approaches have been proposed for network intrusion
detection, where individual agents monitor distinct network segments or protocol layers; for
threat intelligence aggregation, where agents process feeds from different sources; and for
adaptive response systems, where specialised agents handle different stages of the
detection-containment-remediation pipeline \cite{MultiAgentAdaptiveDetection2023,
AIenhancedDefense2022}.

The case for decomposing ransomware detection into modality-specialised agents rests on a
fundamental observation: the statistical properties of static PE-format features and dynamic
runtime API traces are sufficiently different that a single unified model trained on a
concatenation of both modalities is unlikely to learn optimally adapted representations for
each \cite{Ferdous2024Comp}.
Static features are high-dimensional and sparse, with strong category-level interactions;
dynamic features are sequential and exhibit temporal autocorrelation within a single execution
trace.
Modality-specialised agents trained independently on each signal type can develop
representations matched to the distributional properties of their respective inputs, with a
fusion layer responsible for combining the resulting predictions \cite{MultiAgentAdaptiveDetection2023}.

% ------------------------------------------------------------
\subsection{Agent Fusion, Triage, and Escalation}
\label{sec:lit_fusion}
% ------------------------------------------------------------

The problem of combining the predictions of multiple specialised classifiers — \emph{ensemble
fusion} — has a substantial literature in pattern recognition and machine
learning \cite{Ferdous2024Comp}.
Simple fusion strategies include majority voting, averaging of posterior probabilities, and
weighted averaging with weights proportional to individual agent accuracy.
Learned fusion, in which a meta-learner (often a small neural network or logistic regression
model) is trained on the outputs of the individual classifiers, generally outperforms fixed
heuristic fusion by adapting the combination weights to the joint distribution of agent
predictions on the training data \cite{MultiAgentAdaptiveDetection2023}.

An important dimension not adequately addressed in the existing fusion literature is the role
of \emph{prediction uncertainty} in the fusion process.
Conventional fusion methods combine agent probability estimates as if they were equally
reliable across all inputs; in practice, individual agents may be highly confident on inputs
within their training distribution and substantially uncertain on novel or out-of-distribution
inputs.
A fusion scheme that ignores this variation treats a confident prediction and an uncertain
prediction with equal weight, potentially allowing a highly uncertain agent to dominate the
fused output in exactly those cases where its prediction is least reliable.
The inverse-uncertainty-weighted Peer Consultation Protocol developed in Chapter~5 of this
thesis addresses this gap directly, adjusting inter-agent fusion weights at inference time as
a function of each agent's current epistemic uncertainty.

The concept of \emph{triage} — routing detection outputs to different response pathways based
on the system's confidence — is well established in clinical decision support literature but
has received limited systematic treatment in the cybersecurity machine learning context.
Farzadpour et al.\ \cite{Farzadpour2024Thresholding} discuss threshold selection for
security classifiers but focus exclusively on discriminative accuracy rather than on the
operational trade-off between autonomous action and human escalation.
Jointly optimising triage thresholds to maximise detection sensitivity while minimising
unnecessary escalation — particularly using Bayesian optimisation over the combined
objective — represents a methodological gap directly addressed in Chapter~4.

% ============================================================
\section{Knowledge-Augmented and Neuro-Symbolic Detection}
\label{sec:lit_kbs}
% ============================================================

% ------------------------------------------------------------
\subsection{Neuro-Symbolic Artificial Intelligence}
\label{sec:lit_neurosym}
% ------------------------------------------------------------

Neuro-symbolic AI refers to architectures and frameworks that integrate neural learning
components — capable of acquiring representations from raw data — with symbolic reasoning
components that encode structured knowledge, logical constraints, or semantic
relationships \cite{Besold2021neurosymbolicsurvey}.
The motivation for such integration is well-established: neural networks generalise poorly
to novel combinations of concepts not seen during training (the compositional generalisation
problem), and they cannot readily incorporate domain knowledge expressed in symbolic form,
such as logical rules, ontologies, or expert-curated decision criteria.
Conversely, symbolic systems struggle with the uncertainty and noise inherent in real-world
sensor data but can reason reliably over cleanly structured symbolic representations.

In the context of malware detection, the neuro-symbolic paradigm suggests augmenting
purely data-driven detectors with structured representations of threat-intelligence knowledge:
documented attack patterns, tool associations, MITRE ATT\&CK tactic sequences, and
platform-specific vulnerability profiles \cite{Besold2021neurosymbolicsurvey, tian2025llmsurvey}.
Such knowledge is already documented in practitioner-curated repositories — CVE databases,
MITRE ATT\&CK, threat actor profiles, DFIR investigation reports — but has not previously
been operationalised as an independent agent-level signal within a multi-agent detection
framework.
Chapter~5 of this thesis introduces the Knowledge Agent as the component that bridges this
gap, training a neural classifier on practitioner-sourced threat-intelligence features and
integrating its output alongside the data-driven Semantic and Behavioural Agents.

% ------------------------------------------------------------
\subsection{Threat Intelligence as a Detection Signal}
\label{sec:lit_ti}
% ------------------------------------------------------------

Threat intelligence — structured information about adversary capabilities, intentions, and
tactics derived from observed incidents, practitioner analysis, and automated
collection — constitutes a rich source of detection signal that is systematically under-exploited
in the machine learning-based ransomware detection literature.
The MITRE ATT\&CK framework provides a standardised taxonomy of adversary techniques
documented across hundreds of ransomware and advanced persistent threat campaigns, enabling
quantitative characterisation of attack breadth and complexity \cite{razaulla2023age}.
Studies that have incorporated ATT\&CK-derived features into detection pipelines report that
tactic prevalence and tool association profiles are highly discriminative of ransomware
deployments compared to benign or low-sophistication malicious software \cite{Alraizza2025StorageAP}.

However, the integration of threat intelligence into automated detection systems faces two
significant challenges.
First, threat-intelligence data is typically expressed at a higher level of abstraction than
the raw feature vectors used in ML pipelines: bridging the gap between practitioner
narratives and machine-readable feature representations requires structured annotation
effort \cite{Ispahany2024Survey}.
Second, threat-intelligence sources are heterogeneous in quality and recency, and the mapping
from observed tactic indicators to ground-truth maliciousness labels is inherently
probabilistic.
The RansKnow dataset constructed for Chapter~5 addresses these challenges by grounding
features directly in the MITRE ATT\&CK taxonomy and using practitioner-validated thresholds
for label assignment, providing a principled knowledge-grounded training signal for the
Knowledge Agent.

% ============================================================
\section{Concept Drift and Continual Learning}
\label{sec:lit_cd}
% ============================================================

% ------------------------------------------------------------
\subsection{Concept Drift in Malware Detection}
\label{sec:lit_drift}
% ------------------------------------------------------------

Concept drift occurs when the statistical relationship between input features and class labels
changes over time, causing a classifier trained on a historical data distribution to degrade
in performance on more recently collected
samples \cite{Ceschin2022ConceptDrift, Fernando2024FeSAD}.
In the malware detection domain, concept drift is not a rare or exceptional event but an
endemic and deliberate phenomenon: adversaries actively modify their malware to evade detection
systems, introducing distributional shift into both the feature distributions of malicious
samples and, through false positive reduction pressure, into the distribution of borderline
benign software.

Ceschin et al.\ \cite{Ceschin2022ConceptDrift} provide a systematic empirical analysis of
concept drift in malware classifiers, demonstrating that accuracy degradation is measurable
within months of model deployment and that models trained on older data distributions
generalise poorly to current threats even when evaluated on samples from the same broad family
taxonomy.
Fernando et al.\ \cite{Fernando2024FeSAD} propose semi-supervised adaptation mechanisms that
update model representations continuously as new unlabelled samples arrive, reducing the
frequency of model retraining.
Both studies, however, focus on single-model classifiers without explicit uncertainty
quantification or the ability to route uncertain samples to human analysts for labelling.

% ------------------------------------------------------------
\subsection{Continual Learning and Experience Replay}
\label{sec:lit_cl}
% ============================================================

Continual learning addresses the problem of updating a model's knowledge over time as new
labelled examples become available, without access to the full historical training set and
without catastrophically forgetting previously learned
representations \cite{rolnick2019experience}.
The \emph{catastrophic forgetting} problem — the tendency of a neural network updated on new
data to overwrite parameters encoding earlier knowledge — is the central challenge of
continual learning: gradient updates that minimise loss on the current batch can simultaneously
maximise loss on previously learned patterns if the two distributions are
dissimilar \cite{rolnick2019experience}.

Experience replay addresses this problem by maintaining a buffer of previously observed
examples and including a random sample from this buffer in each gradient update batch alongside
the new examples.
The mixed gradient update adjusts model parameters to simultaneously reduce loss on current
and historical samples, mitigating forgetting without requiring access to the full historical
dataset \cite{rolnick2019experience}.
Reservoir sampling \cite{vitter1985random} provides a space-efficient buffer maintenance
strategy: the buffer maintains a fixed-size uniform random sample of all examples seen to
date, with each new example replacing a randomly selected buffer entry with probability equal
to the ratio of buffer size to examples seen.
This ensures that the buffer is an unbiased representative sample of the entire historical
distribution regardless of its temporal structure.

In the ransomware detection context, the Expert Feedback Loop in Chapter~5 operationalises
experience replay in a specific and practically motivated way: human analyst labels on
escalated cases are fed back into the system as the source of new labelled examples, and the
Decision MLP (the only component updated) is trained on the mixed current-plus-replayed
batch.
This design ensures that analyst corrections to systematic model errors are incorporated
incrementally without destabilising the learned representations of the upstream CNN agents.

% ============================================================
\section{Synthesis: Research Gaps and Thesis Positioning}
\label{sec:lit_gaps}
% ============================================================

The preceding review reveals several persistent limitations in the existing ransomware
detection literature that collectively motivate the three empirical investigations in this
thesis.
These gaps are summarised individually before being tabulated in Table~\ref{tab:gaps}.

\textbf{Gap~1: Absence of operational uncertainty quantification.}
The overwhelming majority of deep learning ransomware detectors produce scalar probability
outputs without any characterisation of the reliability of those estimates on individual
inputs \cite{gawlikowski2023Uncertainity, zahoora2022Deep}.
The resulting systems are poorly calibrated and cannot distinguish between high-confidence
predictions on clear-cut cases and low-confidence predictions on borderline or
out-of-distribution samples.
This prevents the construction of principled triage policies that balance autonomous action
against human escalation based on model confidence.

\textbf{Gap~2: Class-imbalance degradation without score-space augmentation.}
Public ransomware datasets are substantially imbalanced, with benign samples outnumbering
ransomware samples in many collections.
Standard augmentation methods such as SMOTE \cite{chawla2002smote} operate in the raw feature
space and may generate implausible synthetic samples in high-dimensional spaces.
No prior work has applied score-space augmentation to the fused prediction vectors of a
multi-agent system, where the dimensionality is low and the conditional distributions of
agent scores are well-suited to generative modelling.

\textbf{Gap~3: Single-modality or na\"{i}vely fused hybrid detection.}
Studies that combine static and dynamic features typically do so by concatenating feature
vectors into a single input to a single unified model \cite{Ferdous2024Comp}, sacrificing
modality-specific representation learning and producing a single aggregate confidence score
that carries no information about which modality was driving the prediction.
No prior work constructs modality-specialised agents with independent uncertainty estimates
that are fused through a learned decision component with uncertainty-aware weighting.

\textbf{Gap~4: No uncertainty-driven triage with Bayesian threshold optimisation.}
While threshold selection for security classifiers is discussed in the
literature \cite{Farzadpour2024Thresholding}, no prior work jointly optimises the
detection threshold and the uncertainty threshold governing autonomous versus escalated
decisions using Bayesian optimisation over a combined discriminative-operational objective
function.

\textbf{Gap~5: Exclusion of structured threat-intelligence from agent-level detection.}
The MITRE ATT\&CK framework and practitioner-curated threat-intelligence sources encode
substantial knowledge about ransomware deployment patterns that is not captured in
measurement-derived static or dynamic features \cite{razaulla2023age, Alraizza2025StorageAP}.
No prior work operationalises this knowledge as an independent neural agent whose uncertainty
output participates in a multi-agent triage framework alongside data-driven agents.

\textbf{Gap~6: Static agent cooperation without uncertainty communication.}
In existing multi-agent detection systems, agents contribute their outputs to a fusion layer
without any mechanism for communicating their current predictive reliability to one
another \cite{MultiAgentAdaptiveDetection2023}.
A high-uncertainty agent contributes the same weight to the fused prediction as a confident
agent, regardless of distributional shift or feature availability.

\textbf{Gap~7: No continual learning under analyst feedback in multi-agent ransomware detection.}
Continual learning methods have been applied to single-model malware classifiers
\cite{Ceschin2022ConceptDrift, Fernando2024FeSAD}, but no prior work incorporates an
analyst-in-the-loop feedback mechanism into a multi-agent detection framework that updates
the system's decision logic incrementally from escalated cases while preserving the learned
representations of the upstream agent backbone networks.

\begin{table}[ht]
\centering
\caption{Identified research gaps and their correspondence to the empirical contributions of
this thesis. Representative prior works that partially address each gap are listed in
column~2; column~3 identifies the chapter that resolves it.}
\label{tab:gaps}
\renewcommand{\arraystretch}{1.25}
\begin{tabular}{@{}p{5.5cm} p{4.5cm} l@{}}
\toprule
\textbf{Research Gap} & \textbf{Representative Prior Work} & \textbf{Addressed In} \\
\midrule
G1: No operational uncertainty quantification in ransomware detection
    & \cite{zahoora2022Deep, gal2016dropout}
    & Chapter~3 \\[4pt]
G2: Class-imbalance augmentation without score-space GAN methods
    & \cite{chawla2002smote, xu2019ctgan}
    & Chapter~3, 4 \\[4pt]
G3: Single-modality or na\"{i}vely fused hybrid detection
    & \cite{Ferdous2024Comp, yumlembam2024comprehensive}
    & Chapter~4 \\[4pt]
G4: No Bayesian uncertainty-threshold triage optimisation
    & \cite{Farzadpour2024Thresholding, optuna2019}
    & Chapter~4, 5 \\[4pt]
G5: Threat-intelligence excluded from agent-level detection
    & \cite{razaulla2023age, Alraizza2025StorageAP}
    & Chapter~5 \\[4pt]
G6: Static agent cooperation without uncertainty communication
    & \cite{MultiAgentAdaptiveDetection2023, gawlikowski2023Uncertainity}
    & Chapter~5 \\[4pt]
G7: No analyst feedback loop in multi-agent ransomware detection
    & \cite{Ceschin2022ConceptDrift, rolnick2019experience}
    & Chapter~5 \\
\bottomrule
\end{tabular}
\end{table}

Taken together, these gaps define a coherent research programme: the limitations of
single-agent detection (G1, G2) motivate the single-agent uncertainty-aware baseline in
Chapter~3; the limitations of single-modality and unfused hybrid detection (G3, G4) motivate
the two-agent multi-modal framework in Chapter~4; and the absence of knowledge integration,
uncertainty communication, and continual learning (G5, G6, G7) motivate the three-agent
extended architecture in Chapter~5.
Each empirical chapter is therefore not an independent contribution but a principled step in
a cumulative research programme, with each chapter's design directly addressing the
limitations identified in the previous.

% ============================================================
\section{Summary}
\label{sec:lit_summary}
% ============================================================

This chapter has reviewed the body of knowledge underpinning the three empirical investigations
in this thesis.
The ransomware threat landscape is characterised by growing operational and economic impact,
rapid adversarial adaptation, and the deployment of multi-stage attack chains that span
reconnaissance, lateral movement, and encryption payload delivery.
Machine learning-based detection has progressed from classical handcrafted-feature methods
through CNN and recurrent deep learning architectures, achieving strong in-distribution
performance but exhibiting systematic overconfidence and limited multi-modal reasoning.
Uncertainty quantification via Monte Carlo Dropout, calibration via temperature scaling, and
explainability via SHAP and LIME attribution provide the technical foundations for
analyst-interpretable, confidence-aware detection.
Multi-agent architectures and neuro-symbolic knowledge integration offer principled pathways
beyond single-model detection, and continual learning via experience replay addresses the
endemic problem of concept drift in evolving malware populations.
Seven specific research gaps have been identified and mapped to the empirical contributions
of Chapters~3, 4, and~5, establishing the logical necessity and coherent progression of the
investigations that follow.
